problem:  
Let \((X,L)\) be a polarized compact Kähler manifold and \(D\subset X\) a smooth divisor.  For each rational \(0<\beta<1\), let \(\omega_{\beta}^{\mathrm{cone}}\) be a Kähler metric on \(X\setminus D\) with cone angle \(2\pi(1-\beta)\) along \(D\), compatible with the holomorphic structure of \(\mathcal O_X(D)\).  
Let \(H^0_{(2)}(X,L^k; \beta)\) denote the finite‑dimensional Hilbert space of holomorphic sections of \(L^k\) that are holomorphic on \(X\setminus D\) and square‑integrable with respect to the cone measure
\[
\|s\|_{k,\beta}^2
=\int_X |s|_{h^k}^2\, |s_D|^{-2\beta}\, \tfrac{\omega_{\beta}^{\mathrm{cone}\,n}}{n!}.
\]
This defines the *cone Bergman metric*
\[
\omega_{k,\beta}=\tfrac1k \Phi_{k,\beta}^*(\omega_{\mathrm{FS}})
\in c_1(L),
\]
determined by an orthonormal basis of \(H^0_{(2)}(X,L^k; \beta)\).

Let \(S(\omega_{k,\beta})\) denote the scalar curvature of \(\omega_{k,\beta}\) on \(X\setminus D\) and \(S_{\beta}^{\mathrm{avg}}\) its average over \(X\) (using the smooth measure \(|s_D|^{-2\beta}\omega_{\beta}^{\mathrm{cone}\,n}\)).

**Question.**  
1. Does the asymptotic *Calabi variance*  
   \[
   \mathcal C_{k,\beta}:=k^{-1}\!\int_{X\setminus D}\!\!
      (S(\omega_{k,\beta})-S_{\beta}^{\mathrm{avg}})^2
      |s_D|^{-2\beta}\omega_{k,\beta}^{\,n}
   \]
   converge as \(k\to\infty\) to a finite limit
   \(\mathcal C_{\infty,\beta}\in[0,+\infty)\)?  
2. If so, does \(\mathcal C_{\infty,\beta}\) vanish if and only if
   a conic cscK metric of angle \(2\pi(1-\beta)\) exists?  
3. Moreover, can the *zero locus* of \(\beta\mapsto \mathcal C_{\infty,\beta}\)
   coincide with the algebraic log K‑stability region
   \(\{\beta: \delta_\beta(X,D)\ge 1\}\)?  
In other words, do *quantized scalar curvature fluctuations* of cone Bergman metrics analytically recover the algebraic δ‑invariant threshold governing log stability?

Why is it a "good" problem:  
This formulation is now mathematically precise and conceptually deep.  It defines a concrete analytic quantity—the asymptotic normalized \(L^2\) variance of scalar curvature for *cone Bergman metrics*—that is well‑posed within the theory of L² holomorphic sections with cone weights.  It asks whether this differential‑geometric invariant detects the existence of conic cscK metrics and coincides with the algebraic δ‑stability threshold.  Such a result would establish a new *quantized curvature criterion for log K‑stability*, geometrically interpreting δ‑invariants in terms of Bergman kernel asymptotics.  The problem’s strength lies in its simplicity and universality: it potentially unifies analytic energy coercivity, conic pluripotential theory, and non‑Archimedean stability into a single measurable curvature‑variance principle—a narrow entrance with profound depth, making it a genuinely “good” and forward‑looking research problem.